% Section 1 -- Introduction
%
% Shape:
%   1.1 Motivating gap: KGE embeddings no longer fit in GPU memory,
%       and the fix for a single GPU does not generalize to many.
%   1.2 Why multi-GPU COVER-replication is the wrong baseline.
%   1.3 Thesis: the scheduler needs a typed, costed IR.
%   1.4 Contributions + results preview.

Knowledge graph embedding (KGE) training pushes two tensors through
the GPU on every step: a block of embedding rows and a matching block
of optimizer state. Both scale linearly with the number of nodes and
with the embedding dimension $d$, and both must be writable on the
device that touches them. At today's dataset sizes that product no
longer fits. A Twitter-scale graph with 41.6\,M nodes at $d\!=\!400$ in
fp32 requires \SI{66.6}{GB} for the embedding matrix and an equal
\SI{66.6}{GB} for Adagrad's accumulator -- a combined \SI{133}{GB}
against a \SI{24}{GB} device-memory budget. Billion-edge KGE
has been an \emph{out-of-core} workload for several years now, and the
community has responded with partitioned trainers that stream
residency in and out of the GPU on a schedule.

\paragraph{The single-GPU story has stabilized.}
Recent work on partitioned KGE -- GE\textsuperscript{2}~\cite{ge2_report},
Marius~\cite{marius}, DGL-KE~\cite{dgl_ke}, PyTorch-BigGraph~\cite{pytorch_biggraph} --
has converged on a common recipe: partition the node set into $p$
blocks, schedule microsteps over the resulting $p^2$ edge buckets so
that each block is resident when its edges are trained, and keep the
blocks on the GPU long enough to amortize loading. GE\textsuperscript{2}'s
COVER algorithm formalizes this as a Partition Scheduling Problem and
reduces the schedule to a Resolvable Balanced Incomplete Block Design
(RBIBD) with parameters $(p, q, 1)$. The result is optimal in admit
count for single-GPU execution and matches the communication
lower-bound $\tfrac{4GNd}{3}$ to within a small constant factor.

\paragraph{The multi-GPU extension is not a schedule.}
The natural way to use two GPUs with a COVER-based trainer is to run
two independent COVER instances, each on a disjoint subset of rounds.
This is what every current partitioned KGE system actually does.
It scales throughput approximately linearly in the common case, but
it makes two decisions silently:

\begin{enumerate}
\item \emph{Inter-GPU handoffs are charged as host reloads.} If lane
0 held partition $\pi$ in round $r$ and lane 1 needs $\pi$ in round
$r\!+\!1$, the system evicts $\pi$ from lane 0 and re-admits it to
lane 1 from the host -- even though $\pi$'s most recent authoritative
copy lives one PCIe hop away on a peer device. The peer path is
invisible to the scheduler.
\item \emph{Lane assignment is a side effect of the partition
numbering.} Because each lane runs its own COVER, the lane a
partition lands on in round $r$ is determined by the order COVER
enumerates microsteps, not by any attempt to keep a partition on the
same lane across rounds. A favorable assignment happens by accident;
an unfavorable one happens as often.
\end{enumerate}

Both are scheduling decisions, but the multi-GPU system has nowhere to
express them. GE\textsuperscript{2}'s scheduler is a COVER solver; its
output is a list of per-device buckets and a host reload pattern. It
does not carry a notion of a \emph{cross-lane handoff}, so it has no
way to reason about one, optimize against one, or hand one to a
runtime that could execute it peer-to-peer.

\paragraph{Thesis.}
We argue that the multi-GPU gap is a representation problem before it
is a performance problem. The scheduler is not missing an algorithm;
it is missing a type. If a compiler can see that a partition is
resident on lane 0 at round $r$ and needed on lane 1 at round $r\!+\!1$,
it can choose whether to serve that admit from host memory or from
a peer device, weigh the two against each other under an explicit
cost model, and commit to one. Today's systems cannot even formulate
the question, because their output format has no field for the
answer.

Our contribution is \emph{Compiled Stateflow}, a typed intermediate
representation (IR) for out-of-core partition schedules, and a
lane-matching compiler that produces plans in that IR. Stateflow
makes every byte of partition residency, host admit, and cross-device
handoff a first-class value with an explicit cost. A plan in
Stateflow is a data structure -- produced, validated, and scored
offline -- that carries four distinct handoff modes
(\texttt{FULL\_RELOAD}, \texttt{ROTATING\_OVERWRITE},
\texttt{DELAYED\_KEEP\_ALIVE}, \texttt{PEER\_RELAY}) and a
decomposable cost function over admits, bucket balance, boundary
crossings, and lane imbalance. The runtime consumes
\texttt{PeerHandoffDescriptor}s projected from the plan and can
execute them as \texttt{cudaMemcpyPeerAsync} without any scheduling
logic of its own.

\paragraph{Results preview.}
On Twitter (41.6\,M nodes, 1.46\,B edges), a 16-partition configuration
on two \SI{24}{GB} RTX A5000s, our lane-matching compiler produces
plans that:
\begin{itemize}
\item lower the host-denominated selection cost from
$1.499 \!\times\! 10^{11}$ to $1.291 \!\times\! 10^{11}$ -- a
\textbf{14\% reduction} in compile-time cost,
\item cut host admits from 72 to 62 (14\%),
\item \textbf{eliminate every \texttt{FULL\_RELOAD} handoff}
(10 $\rightarrow$ 0), converting them into 18
\texttt{ROTATING\_OVERWRITE} and 18 \texttt{PEER\_RELAY} descriptors,
\item expose \textbf{\SI{37.49}{GB} of planned device-to-device traffic}
for the runtime, versus \SI{16.66}{GB} for
\textsc{Disjoint\_Rounds},
\item and, in the final embeddings-scope runtime, execute
\textbf{all 18 planned peer relays} for \textsc{Lane\_Matched}
(\SI{18.74}{GB}, zero fallback, zero descriptor mismatch), while
leaving epoch time effectively flat
(\SI{209.450}{s} vs.\ \SI{208.256}{s} for the host-reload baseline).
\end{itemize}

The compile-time wins are attributable to the compiler; the runtime
result shows that the projected descriptors are now executable and
measurable end-to-end. On this PCIe machine, Phase 4b is
correctness-positive but not yet wall-clock-positive.

\paragraph{Contributions.}
\begin{enumerate}
\item \textbf{Compiled Stateflow}, a typed IR for out-of-core
partition scheduling that separates intra-lane from cross-lane
handoffs, makes every transfer a typed object, and supports a
decomposable cost model (Section~\ref{sec:ir}).
\item \textbf{A lane-matching compiler} for multi-GPU KGE training
that emits Stateflow plans with cross-lane \texttt{PEER\_RELAY}
descriptors, under a host-denominated cost model that does not
credit itself by default for peer savings that are reported but not
optimized
(Section~\ref{sec:compiler}).
\item \textbf{A five-invariant validator} (V1--V5) that rejects
structurally inconsistent or unexecutable plans before they reach
the runtime (Section~\ref{sec:ir:validator}).
\item \textbf{A peer-aware runtime path} built around
\texttt{cudaMemcpyPeerAsync}, gated behind a kill switch and a
fallback to host reloads when peer access is unavailable
(Section~\ref{sec:runtime}).
\item \textbf{An evaluation} on Twitter (41.6\,M, 1.46\,B) showing
the compiler's compile-time gains and the runtime's exact zero-fallback
execution of planned peer relays; on our PCIe A5000 setup, the
peer-enabled path is effectively tied with the host baseline rather
than faster (Section~\ref{sec:eval}).
\end{enumerate}

The rest of the paper is organized as follows.
Section~\ref{sec:background} reviews KGE training, out-of-core
partitioning, and the COVER baseline.
Section~\ref{sec:ir} introduces the Stateflow IR, its types, cost
model, and validator.
Section~\ref{sec:compiler} presents the lane-matching compiler.
Section~\ref{sec:runtime} describes the peer-aware runtime.
Section~\ref{sec:eval} evaluates both on Twitter.
Section~\ref{sec:related} relates Stateflow to other IRs and other
out-of-core trainers, and Section~\ref{sec:conclusion} concludes.
